\begin{beginnerstatement}[In simple terms: control.py]

\path{tools/control.py} is a Python script that serves as the template's
command center. A terminal is a text-based window in which programs are
controlled through typed input. In the terminal, the script provides a
\emph{command-line interface} (CLI). A command such as \texttt{run} selects a
task; an argument supplies additional information for that task.
\emph{Tooling} refers to helper tools for workflows such as installation,
testing, and building. An exit code is a number through which the script
reports success, a domain error, or incorrect usage. For an unexpected Python
error, a traceback shows the sequence of code locations involved;
expected usage and profile errors are displayed without such a detailed
technical report.

The new \texttt{quality} command coordinates such verification tools. A linter
looks for defined errors and rule violations, while a formatter checks the
written form of code. A rule ID is a short, unique identifier such as
\texttt{CQ001}; it keeps a finding recognizable even if its explanation
changes. The dispatcher accepts the command, while dedicated modules contain
the actual quality logic.

\end{beginnerstatement}
